\section{Introduction}

The standard agent loop pays for a model decision at every tool boundary:
model, tool, model, tool, and so on~\citep{yao2023react}. That flexibility is
valuable while a workflow is novel. At scale, however, mature agents often
revisit the same read-only evidence path with the same data dependencies. In
those settings, repeating the model decision adds latency, tokens, and cost
without necessarily adding useful adaptation.

Removing that decision is not ordinary caching. A tool argument may depend on a
specific earlier observation, a repeated sequence may cross an approval or
write, an apparently read-only tool may hide an undeclared effect, and a
program that replays the same tool outputs may still change the downstream
answer. The problem is therefore not just ``find recurring traces.'' The
problem is to decide when a recurrent model-mediated region can be replaced by a
deterministic program and when it should be refused.

Consider the smallest version of this decision in the benchmark.  A NESTFUL record asks
for the percentage increase from 50 dollars to 60 dollars and records
\code{subtract(60,50) -> divide(result,50) -> multiply(result,100)}.  The trace is
executable and every later argument is grounded in an earlier value, so a recurrence-only
optimizer would have a plausible macro.  The benchmark even supplies successful replay:
across 32 attempted families, 24 held-out windows pass, 12 abstain, and none are wrong.
But after provenance and group filtering the strongest candidate family has only 26
independent groups, while the configured exact gate needs 92 zero-violation groups
($\alpha=.05$, $\delta=.10$, eleven thresholds).  GAC therefore retires every family and
keeps the unchanged agent.  The example makes the distinction concrete: recurrence says
what may be compilable; evidence determines whether replacement is justified.

\begin{figure}[t]
  \centering
  \includegraphics[width=\linewidth]{../generated_figures/gac_aha_example.pdf}
  \Description{A benchmark-grounded arithmetic trace looks safe to compile, but a
  recurrence-and-replay rule would ship while GAC retires because 26 independent groups
  fall short of the 92-group exact-gate requirement.}
  \caption{The recurrence-versus-admission distinction on NESTFUL.  This is provider-free
  structural compiler evidence, not a live planning or production-safety result.}
  \label{fig:aha-example}
\end{figure}

The benchmark also shows why tool replay is not full substitution evidence.  In an
earlier 18-record GitHub study, a more aggressive three-read artifact had 45/45 clean
tool replays but only 17/18 exact downstream answers: one Markdown URL was lost from the
continuation.  A checked continuation renderer recovered 18/18 in provider-free replay,
but the paper treats that as a detection mechanism rather than a live safety claim.

Framed this way, the closest precedent is not prompt optimization but the
tracing just-in-time compiler. Dynamo~\citep{bala2000dynamo}, trace-based type
specialization~\citep{gal2009tracejit}, and meta-tracing~\citep{bolz2009pypy}
all detect a hot path, compile it behind a \emph{guard}, and deoptimize to the
interpreter when the guard fails~\citep{holzle1992deopt}. Agent traces admit
the same structure --- hot-trace detection becomes family mining, guard
synthesis becomes an entry contract, and deoptimization becomes fallback to the
unchanged agent --- but they add three obligations a JIT does not face. Tool
arguments must be shown to be \emph{grounded} in observable state rather than
invented by the model; declared \emph{effects} must permit speculation, because
an agent's ``interpreter'' can call the outside world; and the residual error
rate of a compiled region is statistical, so it needs a finite-sample bound
rather than a soundness proof.

We study these obligations through \emph{guarded agentic compaction} (\method),
a compile-or-retire trace-to-program system for recurrent read-only prefixes.
The method normalizes executions into a typed episode representation,
reconstructs value provenance for tool arguments, mines recurrent prefix
families, synthesizes bounded programs, induces runtime guards and verifiers,
and admits an artifact only when a finite-sample selective-risk calculation
supports it. Otherwise the baseline agent is retained unchanged
(\cref{fig:pipeline}). The core contribution is thus not ``agent compilation''
in general, but an admissibility argument for trace-derived specialization.

The evidence is deliberately multi-sided. On three distinct GitHub workflow
families with live provider calls, compiled programs preserve or improve exact
held-out task outcomes while reducing requests, tokens, latency, and cost; a
narrower time-forward task then tests transfer across repositories and retains
both success and compile-time retirement. NESTFUL and API-Bank show the
opposite regime, in which recurrent executable traces still fail to justify
deployment. We also retain the strongest practical counterpoint: when a
workflow is already known, hand-written programs can be competitive or tied, so
any claim of automatic superiority would be incorrect.

The paper makes three contributions:
\begin{enumerate}[leftmargin=1.2em,itemsep=1pt,topsep=2pt]
  \item A guarded specialization pipeline (\cref{alg:compile}) that combines
  typed provenance, effect and position barriers, bounded synthesis, runtime
  verification, and finite-sample compile-or-retire admission
  (\cref{alg:calibrate}, \cref{prop:admission}).
  \item A real-provider evaluation across three public-record workflow families,
  plus a narrower cross-repository, time-forward extension that keeps both
  successful artifacts and principled retirements.
  \item A negative-result boundary showing that recurrence and replayability do
  not themselves license deployment: NESTFUL and API-Bank retire under the same
  admission logic, and manual programs constrain any runtime-dominance claim.
\end{enumerate}
